\section{Guarantees}\label{sec:theory}

The results below cover selection and the optional modules. The gate, the weight floor, the near-twin rule and the learned lift are heuristics with no guarantee: the floor drops items after selection, so the bound of this section applies to the set selected and not to the set injected. Section~\ref{sec:experiments} measures them. Proofs are in Appendix~\ref{app:proofs}.

\subsection{Selection}
\begin{proposition}[Selection guarantee]\label{prop:selection}
Let $F$ be the additive objective \eqref{eq:objective} over the candidates with $w\ge0$, let every candidate satisfy $c(m)\le B$, and let $\mathrm{OPT}$ be the best value of $F$ under the budget. The better of density greedy and the best single item has $F(S)\ge\frac12\mathrm{OPT}$.
\end{proposition}

An exact knapsack solver would reach $\mathrm{OPT}$, and we keep the simpler rule because on development data the two are level (Section~\ref{sec:dev}). The library can add a term for covering query facets and a facility-location term against near-duplicates. The objective is then monotone submodular and the classical cost-benefit argument gives $\frac12(1-\frac1e)\approx0.316$ \cite{krause2005note,leskovec2007cost}, with better constants at higher cost \cite{tang2021revisiting,sviridenko2004note,feldman2023practical}. Both terms are off by default, so that bound plays no part in the results below.

\subsection{Selection with estimated probabilities}
The model only estimates the weights, which raises the question of how estimation error reaches the selection.

\begin{proposition}[From probability error to selection loss]\label{prop:calibration}
Let $p_m$ be the true weight of candidate $m$ and $\hat p_m\ge0$ an estimate, and write $e(S)=\sum_{m\in S}|p_m-\hat p_m|$. Let $S^\star$ maximise $\sum_{m\in S}p_m$ under the budget. If $\hat S$ is feasible and $\sum_{m\in\hat S}\hat p_m\ge\alpha\max_{c(S)\le B}\sum_{m\in S}\hat p_m$, then
\[\sum_{m\in\hat S}p_m\ \ge\ \alpha\sum_{m\in S^\star}p_m-\alpha\,e(S^\star)-e(\hat S).\]
Proposition~\ref{prop:selection} gives $\alpha=\frac12$ for the shipped rule. An exact knapsack has $\alpha=1$, and then the loss is at most $e(S^\star\triangle\hat S)$.
\end{proposition}

Only the error on selected and optimal items matters, and not the error over the whole pool, which is why calibration near the top of the ranking counts for more than a good fit on the long tail of candidates. In the shipped system $\hat p$ is the weight $w$ of Algorithm~\ref{alg:retrieve}.

\subsection{Why validity needs its own signal}
\begin{proposition}[Validity-blind learners]\label{prop:blind}
For any $n\ge2$ there is an instance in which $n$ items per round have identical features, exactly one of them is valid, and $B$ fits one item. A learner that observes neither age nor version information earns $1/n$ per round in expectation against an optimum of $1$, and its regret is $(1-\frac1n)T$ at every horizon. A learner that sees the version chain has regret $0$ on the same instance.
\end{proposition}

No amount of feedback helps here, because feedback carries no information about which copy is current. This is the formal reason to keep version chains apart from usefulness.

\subsection{Optional modules}
Consolidation and forgetting are opt-in. The two results below say what they guarantee, and Section~\ref{sec:optional} measures what they cost.

Let $X$ be a cluster of sentences with weights $p$ ($\sum_x p_x=1$), costs $c$ and cosine similarities $s\in[0,1]$ with $s(x,x)=1$, and let $\Phi(Z)=\sum_xp_x\max_{z\in Z}s(x,z)$ and $D(Z)=1-\Phi(Z)$ be the distortion of a summary $Z\subseteq X$.

\begin{theorem}[Rate--distortion for extractive consolidation]\label{thm:rd}
Let $Z_\ell$ be the first $\ell$ picks of density greedy on $\Phi$, before any pick is skipped for lack of room, with rate $R_\ell=c(Z_\ell)$. For every summary $Z^\star$ of rate $R^\star$,
\[D(Z_\ell)\le e^{-R_\ell/R^\star}+\big(1-e^{-R_\ell/R^\star}\big)D(Z^\star).\]
\end{theorem}

Distortion therefore decays exponentially in rate against any reference summary. The bound averages over $p$ and says nothing about a rare needle in the cluster, which is why consolidation here keeps the sources: if a summary is retrieved, \textsc{Expand} recovers any source at the cost of its tokens.

\begin{theorem}[Store stability]\label{thm:store}
Let items arrive as a Poisson process of rate $a$, and let each be accessed by an independent Poisson process of rate $r$, drawn i.i.d.\ from a distribution with $\mathbb{E}[(e^{r\tau}-1)/r]<\infty$. An item is archived after an idle gap of $\tau$. Then an item's active lifetime has mean $(e^{r\tau}-1)/r$ given $r$, and the number of active items converges in distribution to $\mathrm{Poisson}\big(a\,\mathbb{E}[(e^{r\tau}-1)/r]\big)$.
\end{theorem}

Without forgetting the store grows linearly in $t$, and with an idle rule it stays bounded in distribution, biased exponentially toward frequently used items. The theorem assumes independent access processes, whereas real access is bursty.
